\section{CI und Systemtests}
\label{chap:ci-system-testing}

Die Testpyramide verbindet Einheiten, Adapterverträge, Integrationsfälle, Copy-Fixtures
und Dokumentationsprüfungen. Die Copy-Matrix trennt Linux-, Windows- und macOS-Ergebnisse
als noch zu erfassende \Measured{} Evidenz von vorhandenen, plattformneutralen Tests.
Synthetische Fixtures werden dabei nicht als reale Kundenprojekte ausgegeben.

Abbildung~\ref{fig:ci-test-pyramid} ordnet die Testebenen ein; Abbildung~\ref{fig:os-fixture-matrix}
zeigt die ausstehenden Betriebssystem- und Fixture-Nachweise.

\CaseDiagram{ci-test-pyramid}{CI-Testpyramide vom Modulvertrag bis zur Copy-Fixture.}{fig:ci-test-pyramid}
\CaseDiagram{os-fixture-matrix}{Betriebssystem- und Fixture-Matrix als auszuführender Nachweisplan.}{fig:os-fixture-matrix}
\CaseTable{
\begin{tabularx}{0.94\textwidth}{@{}lX@{}}
\toprule
Ebene & Zweck \\
\midrule
Unit & Deterministische Kern- und Adapterverträge \\
Integration & Planung, Transaktion, Migration und Rollback \\
Copy-Fixture & Portabler Austausch ohne Produktverlust \\
Dokumentation & Links, Quellenstruktur und Sprachparität \\
\bottomrule
\end{tabularx}
}{Testebenen der Fallstudie.}{tab:test-matrix}
